% =====================================================================
% Clase -- Geometría 11° -- Trimestre I -- Semana 04
% Sesión 004 (vie 25 sep., 13:00, 40 min)
% Tema: Coordenadas cartesianas · Subtema: Punto medio y división de un segmento
% Hilo conductor del trimestre I: «LORAN: encontrar un barco con dos hipérbolas»
% Tema 1 de la guía: Coordenadas cartesianas (tema:coordenadas)
% Material del docente (no se entrega a los estudiantes).
% =====================================================================
\documentclass[11pt]{article}
\makeatletter\def\input@path{{../../../../../plantillas/estilo/}}\makeatother
\usepackage{mmcantillo}
\guiadelestudiante{../../guia-didactica/guia-periodo-I-hiperbola-navegacion}

\tipodocumento{Clase}
\asignatura{Geometría}
\titulo{Semana 4: la mitad del camino}
\grado{11°}
\periodo{I}

% DBA: matematicas grado 11 · DBA 6 — Modela objetos geométricos en diversos sistemas de
% coordenadas (cartesiano, polar, esférico) y realiza comparaciones y toma decisiones con
% respecto a los modelos. Evidencia de esta sesión: «Reconoce y utiliza distintos sistemas de
% coordenadas para modelar»; el punto medio y la división de un segmento son la herramienta
% con la que en el tema 2 se define la mediatriz como lugar geométrico.

\begin{document}

\encabezado*

\begin{center}
{\renewcommand{\arraystretch}{1.3}
\begin{tabularx}{\linewidth}{@{}l l l X l@{}}
  \toprule
  \textbf{Sesión} & \textbf{Fecha} & \textbf{Min} & \textbf{Subtema} & \textbf{Entregables}\\
  \midrule
  004 & vie 25 sep., 13:00 & 40 & Punto medio y división de un segmento & Se deja tarea\\
  \bottomrule
\end{tabularx}}
\end{center}

Segundo paso del tema 1. La sesión anterior dio la distancia entre dos puntos; esta da el
punto que está \emph{en la mitad} — y, de paso, el que está a un tercio, o a cualquier
fracción del camino. Es lo que necesita un barco para saber dónde va, y lo que necesitaremos
en la sesión 005 para definir la mediatriz sin dibujarla.
Referencia: \refguia{tema:coordenadas}.

% ---------------------------------------------------------------------
\seccion{Sesión 004 — viernes 25 de septiembre · 13:00 · 40 min}

\textbf{Objetivo.} Calcular el punto medio de un segmento, hallar un extremo conocido el punto
medio, y dividir un segmento en partes iguales o en una fracción dada.

\textbf{Materiales.} Tablero cuadriculado; regla; la guía.

\begin{center}
{\renewcommand{\arraystretch}{1.3}
\begin{tabularx}{\linewidth}{@{}l l X@{}}
  \toprule
  \textbf{Momento} & \textbf{Min} & \textbf{Actividad}\\
  \midrule
  Enganche & 0--5 & Dos puntos en el tablero y la pregunta: «¿cuál es el punto que está
    justo en la mitad?» Casi todos lo señalan bien a ojo. «¿Y cómo lo escribirían con
    números, sin mirar el dibujo?»\\
  El promedio & 5--16 & La fórmula es un promedio, y conviene decirlo así: la $x$ del punto
    medio es el promedio de las dos $x$, y lo mismo con la $y$. $M = \left(\frac{x_1+x_2}{2},
    \frac{y_1+y_2}{2}\right)$. Comprobación de sentido común: $M$ siempre queda
    \emph{entre} los dos extremos.\\
  Al revés & 16--26 & Dado un extremo y el punto medio, hallar el otro extremo: se despeja,
    $x_2 = 2x_M - x_1$. Ejemplo en el tablero. Este es el que más cuesta y es la pregunta 1
    de la tarea.\\
  Fracciones del camino & 26--36 & Partir un segmento en cuatro (dos veces el punto medio)
    y el caso general: el punto a una fracción $t$ del camino de $A$ a $B$ es
    $A + t\,(B - A)$. Ejemplo con el barco: de $(0,0)$ a $(12,9)$, un tercio del camino es
    $(4, 3)$, y son 5 de las 15 millas.\\
  Cierre & 36--40 & Dejar la tarea. Anuncio: la sesión 005 usa el punto medio y la distancia
    juntos, para definir la mediatriz como \emph{lugar geométrico} — todos los puntos que
    equidistan de dos dados.\\
  \bottomrule
\end{tabularx}}
\end{center}

\begin{nota}
\textbf{La comprobación de un segundo.} El punto medio siempre queda entre los dos extremos,
en las dos coordenadas. Si a alguien le sale una ordenada mayor que las dos originales (o menor
que las dos), se equivocó en un signo. Es exactamente el caso de la pregunta 4 de la tarea, y
sirve para todo el año.
\end{nota}

\begin{nota}
\textbf{40 minutos son justos.} Si el «al revés» cuesta más de lo previsto, las fracciones del
camino se pueden dejar solo enunciadas —con el ejemplo del barco— y trabajarlas al inicio de la
sesión 005: la mediatriz necesita el punto medio, no la división en tercios.
\end{nota}

\end{document}
